%%%%%%%%%%%%%%%%%%%%%%%%%%%%%%%%%%%%%%%%%%%%%%%%%%%%%%%%%%
%% HEIG-VD
%%
%% Modèle Latex pour laboratoire
%% 
%% Fichiers à changer : 
%%            - constantes.tex
%%            - Tous les fichiers du corps du rapport 
%%
%% Auteur : TTR
%% Date   : 24.03.2020
%%%%%%%%%%%%%%%%%%%%%%%%%%%%%%%%%%%%%%%%%%%%%%%%%%%%%%%%%%

% Paramétrage du papier utilisé
\documentclass[a4paper,12pt,oneside]{report}

% Ne pas toucher, nécessaire à la mise en forme
\usepackage[utf8x]{inputenc} 
\usepackage[T1]{fontenc} 		% Nouvelle norme pour codage des caractères
\usepackage{titlesec, blindtext, color}
\newcommand{\hsp}{\hspace{20pt}}
\titleformat{\chapter}[hang]{\Huge\bfseries}{\thechapter \hsp}{0pt}{\Huge\bfseries}
\titlespacing*{\chapter}{0pt}{-10pt}{\baselineskip}% see titlesec manual
\usepackage{lmodern} 			% Nouvelle forme de la fonte ComputerModern
\usepackage[english,french]{babel} 		% Règles typographiques françaises, césures
\frenchbsetup{GlobalLayoutFrench=false}
\usepackage{gensymb}  			% Afficher le ° $^{\circ}$
\usepackage{graphicx} 			% Insertion images
\usepackage{epstopdf} 			% Convertir les images eps vers pdf
\epstopdfsetup{outdir=./Images/converted_to_pdf/} % Dossier de sortie des images converties
\pagenumbering{arabic}
\usepackage[top=3cm, bottom=3cm, left=3cm, right=3cm, showframe=false]{geometry}
\usepackage{amsmath}
\usepackage{mathrsfs}
\usepackage{color}
\usepackage{lscape}
\usepackage{tocloft}
\usepackage{float}
\usepackage{textcomp}
\usepackage{markdown}
\usepackage{tabto}
\usepackage{makecell}
\usepackage{pdfpages}
\usepackage{hyperref}
\usepackage{xurl}
\usepackage{xcolor}
\usepackage{tabularx}
\renewcommand\tabularxcolumn[1]{m{#1}}% for vertical centering text in X column
\usepackage{siunitx}
\usepackage{subfiles}
\usepackage{relsize}
\usepackage{pdflscape}

\usepackage{rotating}
\usepackage{tikz}
\usetikzlibrary{shapes,arrows}
\usetikzlibrary{trees}

\usepackage{caption}
\usepackage{subcaption}

\usepackage{arydshln}
\usepackage{blindtext}

\usepackage{afterpage}
\usepackage{bbding}

\usepackage{multirow}
\newcolumntype{M}[1]{>{\centering\arraybackslash}m{#1}}
%\usepackage{verbatim}

\usepackage{amssymb}% http://ctan.org/pkg/amssymb
\usepackage{pifont}% http://ctan.org/pkg/pifont
\newcommand{\cmark}{\ding{51}}%
\newcommand{\xmark}{\ding{55}}%

% Enumerate style
\renewcommand{\labelenumii}{\arabic{enumi}.\arabic{enumii}}
\renewcommand{\labelenumiii}{\arabic{enumi}.\arabic{enumii}.\arabic{enumiii}}
\renewcommand{\labelenumiv}{\arabic{enumi}.\arabic{enumii}.\arabic{enumiii}.\arabic{enumiv}}

% Directory tree package
\usepackage{dirtree}

\definecolor{mygray}{gray}{0.95}

\newcommand{\tree}[4]{
\begin{figure}[H] 
\renewcommand\figurename{Tree}
\centering 
\setlength{\fboxsep}{0pt}%
\setlength{\fboxrule}{0pt}%
\fbox{\colorbox{mygray}{
\begin{minipage}[t]{#1\linewidth}
 \dirtree{#4}	 
 \end{minipage}
 }}
 \caption{#3}
 \label{tree:#2}
\end{figure}}

\definecolor{myOrange}{rgb}{1,0.93,0.86}

\newcommand{\warning}[1]{%
\begin{center}
\setlength{\fboxsep}{4pt}%
\setlength{\fboxrule}{0pt}%
\fbox{\colorbox{myOrange}{\parbox{0.95\linewidth}{\centering{#1}}}}
\end{center}
}

\definecolor{myblue}{rgb}{0.8,0.9,1}

\newcommand{\link}[1]{%
\begin{center}
\setlength{\fboxsep}{4pt}%
\setlength{\fboxrule}{0pt}%
\fbox{\colorbox{myblue}{\parbox{0.95\linewidth}{\centering \url{#1}}}}
\end{center}
}

% Sorted List

\usepackage{datatool}% http://ctan.org/pkg/datatool
\newcommand{\sortitem}[1]{%
  \DTLnewrow{list}% Create a new entry
  \DTLnewdbentry{list}{description}{#1}% Add entry as description
}
\newenvironment{sortedlist}{%
  \DTLifdbexists{list}{\DTLcleardb{list}}{\DTLnewdb{list}}% Create new/discard old list
}{\DTLsort{description}{list}% Sort list
  \begin{itemize}%
    \DTLforeach*{list}{\theDesc=description}{%
      \item \theDesc}% Print each item
  \end{itemize}%
}

% Code in one line\
\usepackage{xcolor}
\usepackage{listings}
\usepackage{xparse}
% Bash listing style

\lstset{literate=
  {á}{{\'a}}1 {é}{{\'e}}1 {í}{{\'i}}1 {ó}{{\'o}}1 {ú}{{\'u}}1
  {Á}{{\'A}}1 {É}{{\'E}}1 {Í}{{\'I}}1 {Ó}{{\'O}}1 {Ú}{{\'U}}1
  {à}{{\`a}}1 {è}{{\`e}}1 {ì}{{\`i}}1 {ò}{{\`o}}1 {ù}{{\`u}}1
  {À}{{\`A}}1 {È}{{\'E}}1 {Ì}{{\`I}}1 {Ò}{{\`O}}1 {Ù}{{\`U}}1
  {ä}{{\"a}}1 {ë}{{\"e}}1 {ï}{{\"i}}1 {ö}{{\"o}}1 {ü}{{\"u}}1
  {Ä}{{\"A}}1 {Ë}{{\"E}}1 {Ï}{{\"I}}1 {Ö}{{\"O}}1 {Ü}{{\"U}}1
  {â}{{\^a}}1 {ê}{{\^e}}1 {î}{{\^i}}1 {ô}{{\^o}}1 {û}{{\^u}}1
  {Â}{{\^A}}1 {Ê}{{\^E}}1 {Î}{{\^I}}1 {Ô}{{\^O}}1 {Û}{{\^U}}1
  {Ã}{{\~A}}1 {ã}{{\~a}}1 {Õ}{{\~O}}1 {õ}{{\~o}}1
  {œ}{{\oe}}1 {Œ}{{\OE}}1 {æ}{{\ae}}1 {Æ}{{\AE}}1 {ß}{{\ss}}1
  {ű}{{\H{u}}}1 {Ű}{{\H{U}}}1 {ő}{{\H{o}}}1 {Ő}{{\H{O}}}1
  {ç}{{\c c}}1 {Ç}{{\c C}}1 {ø}{{\o}}1 {å}{{\r a}}1 {Å}{{\r A}}1
  {€}{{\euro}}1 {£}{{\pounds}}1 {«}{{\guillemotleft}}1
  {»}{{\guillemotright}}1 {ñ}{{\~n}}1 {Ñ}{{\~N}}1 {¿}{{?`}}1
}


\lstdefinestyle{C}{
	backgroundcolor=\color{mygray},   		% choose the background color; you must add \usepackage{color} or \usepackage{xcolor}; should come as last argument
	basicstyle=\footnotesize\ttfamily,  	% the size of the fonts that arexy used for the code
	belowcaptionskip=1\baselineskip,
	breakatwhitespace=true,         		% sets if automatic breaks should only happen at whitespace
	breaklines=true,                 		% sets automatic line breaking
	postbreak=\raisebox{0ex}[0ex][0ex]{$\hookrightarrow\;$},
	captionpos=b,                    		% sets the caption-position to bottom
	commentstyle=\color{green!40!black},,   % comment style
	deletekeywords={...},            		% if you want to delete keywords from the given language
	escapeinside={\%*}{*)},          		% if you want to add LaTeX within your code
	extendedchars=true,              		% lets you use non-ASCII characters; for 8-bits encodings only, does not work with UTF-8
	firstnumber=1,                			% start line enumeration with line 1
	frame=leftline,	                   
	identifierstyle=\color{blue!50!black},
	keepspaces=true,                 		% keeps spaces in text, useful for keeping indentation of code (possibly needs columns=flexible)
	keywordstyle=\bfseries\color{violet},   % keyword style
	language=C,                 			% the language of the code
	morekeywords={*,...},            		% if you want to add more keywords to the set
	numbers=left,                    		% where to put the line-numbers; possible values are (none, left, right)
	numbersep=10pt,                   		% how far the line-numbers are from the code
	numberstyle=\tiny\color{gray}, 			% the style that is used for the line-numbers
	rulecolor=\color{black},         		% if not set, the frame-color may be changed on line-breaks within not-black text (e.g. comments (green here))
	showspaces=false,                		% show spaces everywhere adding particular underscores; it overrides 'showstringspaces'
	showstringspaces=false,          		% underline spaces within strings only
	showtabs=false,                  		% show tabs within strings adding particular underscores
	stepnumber=1,                    		% the step between two line-numbers. If it's 1, each line will be numbered
	stringstyle=\color{orange},
	tabsize=2,	                   			% sets default tabsize to 2 spaces
	title=\lstname,
	%xleftmargin=1.5cm,
	 xleftmargin=\parindent,
}

\definecolor{mygreen}{rgb}{0.1,0.4,0.1}
\lstdefinestyle{console}{
	backgroundcolor=\color{mygray},
	basicstyle=\footnotesize\ttfamily\color{black},
	commentstyle=\color{mygreen},
	keywordstyle=\color{blue},
    numberstyle=\tiny\color{gray},
    stringstyle=\color{purple},
    breakatwhitespace=false,         
    breaklines=true,                 
    captionpos=b,                    
    keepspaces=true,                 
    numbers=left,                    
    numbersep=5pt,                  
    showspaces=false,                
    showstringspaces=false,
    showtabs=false,                  
    tabsize=2,
    language=bash,
    xleftmargin=\parindent,
}
%\lstset{language=Bash,keywordstyle={\bfseries \color{black}}}
%\NewDocumentCommand{\codeword}{v}{\texttt{\colorbox{lightgray}{\textcolor{black}{#1}}}}

\NewDocumentCommand{\codeword}{v}{\textcolor{blue}{#1}}

%-------------------------------------------------------------------------------------

\newcommand{\codetxt}[1]{\textcolor{blue}{\texttt{\path{#1}}}}

%-------------------------------------------------------------------------------------

\newcommand{\eme}{\ieme{ }}
\newcommand{\er}{\ier{ }}
\newcommand{\ere}{\iere{ }}

\usepackage{float,graphicx}

\newcommand{\blue}[1]{\textcolor{blue}{#1}}

%Figure
\newcommand{\figc}[3]{%
  \begin{figure}[H]
    \centering
    \includegraphics[width=#1\textwidth]{Include/Figure/#2.png}
    \caption{#3}
    \label{fig:#2}
  \end{figure}}
  
  \newcommand{\figpdf}[3]{%
  \begin{figure}[H]
    \centering
    \includegraphics[width=#1\textwidth]{Include/Figure/#2.pdf}
    \caption{#3}
    \label{fig:#2}
  \end{figure}}
  

\newcommand{\figl}[4]{
\begin{table}[H]
\begin{tabularx}{\textwidth}{l X}
    \centered{\includegraphics[width=#1\textwidth]{Include/Figure/#2.png}} & \begin{tabular}{#3}#4\end{tabular}
\end{tabularx}
\end{table}}

\newcommand{\figr}[4]{
\begin{table}[H]
\begin{tabularx}{\textwidth}{X r}
	\begin{tabular}{#3} #4 \end{tabular} & \centered{\includegraphics[width=#1\textwidth]{Include/Figure/#2.png}}
\end{tabularx}
\end{table}}

\newcommand{\figcjpg}[3]{%
  \begin{figure}[H]
    \centering
    \includegraphics[width=#1\textwidth]{Include/Figure/#2.jpg}
    \caption{#3}
    \label{fig:#2}
  \end{figure}}
  
\newcommand{\figljpg}[4]{
\begin{table}[H]
\begin{tabularx}{\textwidth}{l X}
    \centered{\includegraphics[width=#1\textwidth]{Include/Figure/#2.jpg}} & \begin{tabular}{#3}#4\end{tabular}
\end{tabularx}
\end{table}}
  
\newcommand{\figrjpg}[4]{
\begin{table}[H]
\begin{tabularx}{\textwidth}{X r}
	\begin{tabular}{#3} #4 \end{tabular} & \centered{\includegraphics[width=#1\textwidth]{Include/Figure/#2.jpg}}
\end{tabularx}
\end{table}}

\newcommand{\centered}[1]{\begin{tabular}{l} #1 \end{tabular}}

\newcommand{\tbox}[1]{
\begin{tabular}{|c|}
\hline #1 \\[4pt]\hline
\end{tabular}}

\newcommand{\ii}[1]{%
\begin{itemize}[\;]%
\item #1%
\end{itemize}}

%TO DO
\newcommand{\todo}[1]{\begin{center}
\colorbox{orange!50}{\Large #1}
\end{center}}

\newcommand{\listequationsname}{\Large{Liste des équations}}
\newlistof{myequations}{equ}{\listequationsname}
\newcommand{\myequations}[1]{%
\addcontentsline{equ}{myequations}{\protect\numberline{\theequation}#1}\par}

\newcommand{\ddl}[1]{\underline{\underline{#1}}}


% fichier avec les constantes (à adapter pour les étudiants)
%%%%%%%  Constantes (à changer) %%%%%%%%

\newcommand{\nomAuteur}{Tristan Traiber}

\newcommand{\titre}{Analog LTE-M / NB-IoT Watch}

\newcommand{\bigTitre}{Master of Science HES-SO in Engineering}
\newcommand{\subBigTitre}{Orientation : Industrial Technologies (TIN)}

\newcommand{\shortTitre}{Analog LTE-M / NB-IoT Watch}

\newcommand{\nameLabo}{Master Thesis Report} % numéro de laboratoire

\newcommand{\professeur}{Medard Rieder}

\newcommand{\assistant}{Patrice Rudaz}

\newcommand{\classe}{HES-SO Valais/Wallis}

\newcommand{\departement}{EIE}

\newcommand{\cours}{ProjApp}

\newcommand{\addr}{Master of Science HES-SO in  Engineering \\ Av. de Provence 6 \\ CH-1007 Lausanne}

% Entête et pieds de page (A ADAPTER)
\pagestyle{empty}
\usepackage{fancyhdr}
\setlength{\headheight}{25pt}
\pagestyle{fancy}
\renewcommand{\chaptermark}[1]{\markboth{#1}{#1}}

% Entête :
\fancypagestyle{plain}{}
\renewcommand\headrulewidth{1pt}
\fancyhead[L]{\includegraphics[height=8mm]{Include/Figure/logo/mse-small-logo.png}}
\fancyhead[R]{\textsc{\nomAuteur} \\ \begin{otherlanguage}{english} \today \end{otherlanguage}}
\fancyhead[C]{\nameLabo \\ \shortTitre}
% Pied de page
\fancyfoot{} % clear all footer fields
\renewcommand{\footrulewidth}{1pt}
\fancyfoot[R]{\thepage} 
\fancyfoot[L]{\textsc{\leftmark}}

\newcommand{\specialcell}[2][l]{%
  \begin{tabular}[#1]{@{}l@{}}#2\end{tabular}}

% Début du corps du rapport 
\begin{document}	

	\pagenumbering{Roman}
	
	% première page
	%%%%%%%%%%%%%%%%%%%%%%%%%%%%%%%%%%%%%%%%%%%%%%%%%%%%%%%%
% Ce document ne doit pas être modifié
% Il faut modifier les informations depuis le fichier
% constantes.tex
%%%%%%%%%%%%%%%%%%%%%%%%%%%%%%%%%%%%%%%%%%%%%%%%%%%%%%%%


\begin{center}
\includegraphics[height=8mm]{Include/Figure/logo/mse-full-cropped.pdf}
\hfill
\includegraphics[height=8mm]{Include/Figure/logo/HES-SO_logo_Pantone.pdf}
\end{center}

\vfill

\begin{flushleft}
	\small \addr
\end{flushleft}

\vfill

\begin{center}
	\huge \bigTitre \\
	\Large \subBigTitre \\
\end{center}

\vfill

\begin{center}
	\Large \nameLabo:\\
	\huge \titre \\
\end{center}

\vfill

\vfill

\begin{flushright}
	\Large Made by\\
	\huge \nomAuteur\\
	
	\vfill
	
	\Large Under the direction of \\
	\normalsize Prof. \professeur\\
	HES-SO Valais/Wallis \\
	Rue de l'Industrie 23 \\
	1950 Sion \\
\end{flushright}
\vfill

\begin{flushright}
	%External expert -\\
\end{flushright}

\vfill

\begin{otherlanguage}{english}
\begin{flushright}
	Lausanne, HES-SO, Master MSE, \today
\end{flushright}
\end{otherlanguage}

	\thispagestyle{empty}

%%%%%%%%%%%%%%%%%%%%%%%%%%%%%%%%%%%%%%%%%%%%%%%%%%%%%%%%%%%%%%%%%%%%%%%%%%%%%%%

	\afterpage{\null\newpage}

%%%%%%%%%%%%%%%%%%%%%%%%%%%%%%%%%%%%%%%%%%%%%%%%%%%%%%%%%%%%%%%%%%%%%%%%%%%%%%%

\newpage

\begin{flushleft}
	Accepted by the HES SO Master MSE (Suisse, Lausanne) on the proposal of \\[12pt]
	Prof. \professeur, TM Project Advisor\\
	Ing. \assistant, Senior academic associate\\[48pt]
\end{flushleft}

\begin{otherlanguage}{english}
\begin{flushleft}
	Lausanne, \today \\[48pt]
\end{flushleft}
\end{otherlanguage}

\begin{flushleft}
	Prof. \professeur \hfill Prof. Philippe Barrade \\
	Project Adviser  \hfill Filliar manager \\
\end{flushleft}


	
	
%%%%%%%%%%%%%%%%%%%%%%%%%%%%%%%%%%%%%%%%%%%%%%%%%%%%%%%%%%%%%%%%%%%%%%%%%%%%%%%
	% enlève le numéro sur la première page
	\newpage
	
	\subfile{2_lexic.tex}

	\newpage

	\begin{otherlanguage}{english}
		\chapter*{Table of Contents}
		\addcontentsline{toc}{chapter}{Table of Contents}
		\renewcommand\contentsname{} % the empty name
		\begingroup
			\let\clearpage\relax
			\vspace{-3.5cm} % the removed space. Set as appropriate
			\tableofcontents
		\endgroup
	\end{otherlanguage}

	\newpage	
	
	
	\pagenumbering{arabic}
	\setcounter{page}{1}
	
	\subfile{3_introduction.tex}
		
	\newpage
	
	\subfile{4_research_part.tex}
	
	\newpage
	
	\subfile{4_b_solution_part.tex}
	
	\newpage	
	
	\subfile{4_1_hardware_part.tex}

	\newpage
	
	\subfile{4_2_antenna_part.tex}	
		
	\newpage
	
	\subfile{4_3_software_part.tex}	

	\newpage	
	
	\subfile{5_experimental_part.tex}
	
%	\newpage
	
%	\subfile{6_result.tex}
	
	\newpage
	
	\subfile{7_discussion.tex}
	
	\newpage
	
	\subfile{8_conclusion.tex}
	
	%\newpage 
	
	%\subfile{report_struct.tex}
	
	%Section pour la signature
	\section*{Signatures}
  	\;\\[-20pt]
    %\vspace{2cm}
    
    \tabto{8cm}Lausanne, \today\\
    \tabto{8cm}Tristan Traiber \\
    
    \begin{flushright}
    \begin{figure}[H]
       \tabto{8cm} \includegraphics[width=4.0cm]{Include/Figure/logo/SignatureTR.PNG}
    \end{figure}
    \end{flushright}

\begin{otherlanguage}{english}
	\bibliographystyle{unsrt}
	\bibliography{bibliographie.bib}
\end{otherlanguage}


\begin{otherlanguage}{english}
	\listoffigures
	\listoftables
\end{otherlanguage}    
   
	
	\newpage
	\subfile{9_annexes}
	
	\newpage
	%\section{tableau}
% Inclure un tableau
\begin{table}[H]
\centering
    \begin{tabular}{|c|c|c|c|c|}
    \hline
    11 & 12 & 13 & 15 & 15  \\
    \hline
    21 & 22 & 23 & 24 & 25  \\
    \hline
    31 & 32 & 33 & 34 & 35  \\
    \hline
    41 & 42 & 43 & 44 & 45  \\      
    \hline
    \end{tabular}
    \caption{Légende}
    \label{table:Name}
\end{table}

\section{équation}
% Inclure une équation
\begin{equation}
    \label{eq:Name}
    (a+b)^2 = a^2 + 2 \cdot a \cdot b + b^2 
\end{equation}
\myequations{Titre de l'équation}

\section{figure}
% Inclure une figure
\begin{figure}[H]
    \centering
    \includegraphics[width=0.5\textwidth]{Include/Figure/logo/SignatureTR.png}
    \caption{Légende}
    \label{fig:Name}
\end{figure}

\section{PDF}
% Inclure un fichier PDF
\begin{figure}[H]
    \centering
    \includegraphics[width=0.8\textwidth]{Include/Figure/logo/mse-full-cropped.pdf}
\end{figure}

\section{itemize}
\begin{itemize}
  \item One
  \item Two
\end{itemize}

\section{enumerate}
\begin{enumerate}
    \item
        \subitem
\end{enumerate}

% Commentaire sur plusieurs lignes
\iffalse 
    Texte à mettre en commentaire
\fi

\section{lstlisting}
\begin{lstlisting}[language=c++,caption={Exemple d'un bon programme d'allocation dynamique}]
code
\end{lstlisting}

\begin{lstlisting}[style=c,label={lst:xxx},caption={}]
 code C
\end{lstlisting}

\begin{lstlisting}[style=console,label={lst:xxx},caption={}]
 code bash
\end{lstlisting}

%\begin{lstlisting}[style=txt,label={lst:xxx},caption={}]
% Text
%\end{lstlisting}
	
%Fin du document
\end{document}